\chapter{Two-by-two matrices and M\"obius maps}\label{app:matrices}
\chaptertag{A compact toolkit for continued fractions}

This appendix develops only the matrix algebra used in Chapter~3. The central
idea is simple: a matrix with determinant $\pm1$ can be reversed without ever
leaving the integers.

\section{Multiplication and determinants}

A $2\times2$ matrix acts on a column vector by
\[
 \begin{pmatrix}a&b\\c&d\end{pmatrix}
 \begin{pmatrix}x\\y\end{pmatrix}
 =\begin{pmatrix}ax+by\\cx+dy\end{pmatrix}.
\]
Matrix multiplication is defined so that composition works as expected:
applying $B$ and then $A$ is the same as applying $AB$. Explicitly,
\[
 \begin{pmatrix}a&b\\c&d\end{pmatrix}
 \begin{pmatrix}e&f\\g&h\end{pmatrix}
 =\begin{pmatrix}ae+bg&af+bh\\ce+dg&cf+dh\end{pmatrix}.
\]

The determinant of $A=\begin{psmallmatrix}a&b\\c&d\end{psmallmatrix}$ is
$\det A=ad-bc$. It measures invertibility:
\[
 A^{-1}=\frac1{ad-bc}
 \begin{pmatrix}d&-b\\-c&a\end{pmatrix},
 \qquad\text{provided }ad-bc\ne0.
\]
The identity $\det(AB)=\det(A)\det(B)$ makes determinants especially useful
for long products.

\begin{definition}
The group $\operatorname{GL}_2(\Z)$ consists of the integer matrices with
determinant $\pm1$. Its subgroup $\operatorname{SL}_2(\Z)$ consists of those
with determinant $+1$.
\end{definition}

\begin{proposition}[Integer reversibility]
An integer $2\times2$ matrix has an inverse with integer entries if and only if
its determinant is $\pm1$.
\end{proposition}

\begin{proof}
If $\det A=\pm1$, the displayed inverse formula has integer entries. Conversely,
if both $A$ and $A^{-1}$ have integer entries, then
$1=\det(I)=\det(A)\det(A^{-1})$ is a product of two integers. Hence
$\det(A)=\pm1$.
\end{proof}

\begin{bigidea}
The columns of a matrix in $\operatorname{GL}_2(\Z)$ form an integer basis of
$\Z^2$: every integer vector has a unique expression as an integer linear
combination of those columns. This is the exact algebraic strength supplied by
the determinant identity for consecutive convergents.
\end{bigidea}

\section{Matrices as fractional linear maps}

Associate to an invertible matrix
$A=\begin{psmallmatrix}a&b\\c&d\end{psmallmatrix}$, with $ad-bc\ne0$,
the M\"obius map
\[
  T_A(t)=\frac{at+b}{ct+d}.
\]
A direct substitution shows
\[
  T_A\bigl(T_B(t)\bigr)=T_{AB}(t).
\]
Thus nested fractional expressions become matrix products. Multiplying $A$ by
any nonzero scalar leaves $T_A$ unchanged, but retaining the specific integer
matrix remembers determinant and divisibility information that the function
alone would forget.

The inverse map is $T_{A^{-1}}$, so a change of variable transports fixed
points by conjugation:
\[
 x=T_N(x),\quad y=T_M(x)
 \quad\Longrightarrow\quad y=T_{MNM^{-1}}(y).
\]
Section~\ref{sec:periodic-fixed-points} uses this identity to prove that
adding a finite prefix to a periodic continued fraction preserves
quadratic irrationality. Section~\ref{sec:galois} then proves the precise
criterion for pure periodicity.

When $ct+d\ne0$, the quotient rule gives
\[
  T_A'(t)=\frac{ad-bc}{(ct+d)^2}.
\]
Therefore the determinant tells us whether the map is increasing or decreasing.
This is the M\"obius derivative trick used to locate alternating convergents.

\section{Continued fractions as iterated products}

One partial quotient corresponds to
\[
  D(a)=\begin{pmatrix}a&1\\1&0\end{pmatrix},
  \qquad T_{D(a)}(t)=a+\frac1t.
\]
Consequently,
\[
 T_{D(a_0)\cdots D(a_k)}(t)
 =[a_0;a_1,\dots,a_k,t].
\]
The convergent recurrences are equivalent to the product identity
\[
 D(a_0)\cdots D(a_k)
 =\begin{pmatrix}p_k&p_{k-1}\\q_k&q_{k-1}\end{pmatrix}.
\]
Since every $D(a_i)$ has determinant $-1$, the product has determinant
$(-1)^{k+1}$. Hence
\[
 p_kq_{k-1}-p_{k-1}q_k=(-1)^{k+1}.
\]

This one product packages four facts:
\begin{itemize}
  \item its first column is the current convergent $(p_k,q_k)$;
  \item its second column is the preceding convergent $(p_{k-1},q_{k-1})$;
  \item its determinant says the two columns form an integer basis;
  \item its M\"obius action inserts an arbitrary remaining tail.
\end{itemize}

\begin{example}
For $[4;2,6,7]=415/93$,
\[
 D(4)D(2)D(6)D(7)
 =\begin{pmatrix}415&58\\93&13\end{pmatrix}.
\]
Its inverse is
\[
 \begin{pmatrix}13&-58\\-93&415\end{pmatrix}.
\]
The top row applied to the first column gives
$13(415)-58(93)=1$. In the opposite direction, multiplying a vector by the
inverse recovers its integer coordinates in the basis of consecutive
convergents.
\end{example}

\begin{exercisebox}
\begin{enumerate}
  \item Multiply three general matrices $D(a_0)D(a_1)D(a_2)$ and identify the
  two convergents in its columns. \exlevel{Fluency}
  \item Prove $\det(AB)=\det(A)\det(B)$ directly for $2\times2$ matrices.
  \exlevel{Reasoning}
  \item Find the inverse of
  $\begin{psmallmatrix}17&7\\12&5\end{psmallmatrix}$ and use it to express
  $(1,0)$ and $(0,1)$ as integer combinations of its columns.
  \exlevel{Fluency}
  \item Verify directly that $T_A\circ T_B=T_{AB}$. Explain why multiplying a
  matrix by $-1$ does not change its M\"obius map. \exlevel{Reasoning}
  \item Let $P=D(a_1)\cdots D(a_\ell)$. Show that a purely periodic tail
  $t=[\overline{a_1,\dots,a_\ell}]$ is a fixed point of $T_P$, and derive its
  quadratic equation from the four entries of $P$. \exlevel{Reasoning}
\end{enumerate}
\end{exercisebox}
